\begin{proof}[Proof of \cref{obj:lem:frobenius-center-certificate}]
Throughout \(q=2(m-1)>0\), so that
\[
q+2=2m .
\]

\emph{Step 1 (the center is feasible).}  The point \((1,1,1)\) has nonnegative coordinates and satisfies \(q\cdot1+1+1=q+2=2m\), hence \((1,1,1)\in T_m\). % lean: frobenius_center_certificate

\emph{Step 2 (a bias--variance decomposition of the weighted norm on \(T_m\)).}  Let \((x,y,z_{\mathrm{sp}})\in T_m\).  Expanding the squares,
\[
q(x-1)^2+(y-1)^2+(z_{\mathrm{sp}}-1)^2
=\bigl(qx^2+y^2+z_{\mathrm{sp}}^2\bigr)
-2\bigl(qx+y+z_{\mathrm{sp}}\bigr)+(q+2),
\]
and substituting the simplex equation \(qx+y+z_{\mathrm{sp}}=2m\) together with \(q+2=2m\) gives \(-2(2m)+(q+2)=-(q+2)\).  Hence
\[
qx^2+y^2+z_{\mathrm{sp}}^2
=(q\cdot1^2+1^2+1^2)+q(x-1)^2+(y-1)^2+(z_{\mathrm{sp}}-1)^2 .
\]
% lean: frobenius_center_certificate

\emph{Step 3 (strict minimality of the center for the norm).}  Suppose \((x,y,z_{\mathrm{sp}})\in T_m\) with \((x,y,z_{\mathrm{sp}})\ne(1,1,1)\).  Each of \(q(x-1)^2\), \((y-1)^2\), \((z_{\mathrm{sp}}-1)^2\) is nonnegative, and if their sum vanished then, since \(q>0\), each would vanish and the point would be \((1,1,1)\).  Therefore
\[
0<q(x-1)^2+(y-1)^2+(z_{\mathrm{sp}}-1)^2 ,
\]
so by Step 2, \(q\cdot1^2+1^2+1^2<qx^2+y^2+z_{\mathrm{sp}}^2\), and taking square roots (both arguments being nonnegative and the square root strictly increasing)
\[
\sqrt{q\cdot1^2+1^2+1^2}<\sqrt{qx^2+y^2+z_{\mathrm{sp}}^2}.
\]
This is the Frobenius-center assertion. % lean: frobenius_center_certificate

\emph{Step 4 (necessity of \(c_x/q=c_y\)).}  Assume \(\phi_{r,\kappa}(1,1,1)\le\phi_{r,\kappa}(x,y,z_{\mathrm{sp}})\) for all \((x,y,z_{\mathrm{sp}})\in T_m\), and consider the curve
\[
t\longmapsto(1+t,\,1-qt,\,1).
\]
It is feasible for \(|t|<\min\{1,1/q\}\): the coordinates \(1+t\) and \(1-qt\) are then positive, the third is \(1\), and
\[
q(1+t)+(1-qt)+1=q+2=2m
\]
for every \(t\), so the curve stays in \(T_m\).  Consequently \(t=0\) is a local minimum of
\[
g(t)=\phi_{r,\kappa}(1+t,1-qt,1)
=c_x(1+t)+c_y(1-qt)+c_z
+\kappa\sqrt{q(1+t)^2+(1-qt)^2+1}.
\]
The radicand has derivative
\[
\frac{d}{dt}\Bigl(q(1+t)^2+(1-qt)^2+1\Bigr)\Big|_{t=0}
=2q(1+t)-2q(1-qt)\Big|_{t=0}=2q-2q=0,
\]
and its value at \(t=0\) is \(q+2>0\), so the chain rule makes the derivative of the square-root term vanish at \(t=0\).  Hence \(g'(0)=c_x-qc_y\), and local minimality at an interior point forces
\[
c_x-qc_y=0,
\qquad\text{i.e.}\qquad
\frac{c_x}{q}=c_y .
\]
% lean: frobeniusCenterTangentXY_deriv

\emph{Step 5 (necessity of \(c_y=c_z\)).}  Consider likewise
\[
t\longmapsto(1,\,1+t,\,1-t),
\]
feasible for \(|t|<1\) since then all three coordinates are nonnegative and \(q\cdot1+(1+t)+(1-t)=q+2=2m\).  Writing \(h(t)=\phi_{r,\kappa}(1,1+t,1-t)\), the radicand \(q+(1+t)^2+(1-t)^2\) has derivative \(2(1+t)-2(1-t)\), which vanishes at \(t=0\), and its value there is \(q+2>0\); so again the square-root term contributes nothing to \(h'(0)\) and
\[
h'(0)=c_y-c_z=0,
\qquad\text{i.e.}\qquad
c_y=c_z .
\]
% lean: frobeniusCenterTangentYZ_deriv

\emph{Step 6 (strict optimality under the balance conditions).}  Suppose finally \(\kappa>0\), \(c_x/q=c_y\) and \(c_y=c_z\); thus \(c_x=qc_y\) and \(c_z=c_y\).  For every \((x,y,z_{\mathrm{sp}})\in T_m\) the linear part is then constant:
\[
c_xx+c_yy+c_zz_{\mathrm{sp}}
=c_y\bigl(qx+y+z_{\mathrm{sp}}\bigr)
=c_y\cdot 2m
=c_y(q+2)
=c_x\cdot1+c_y\cdot1+c_z\cdot1 .
\]
For \((x,y,z_{\mathrm{sp}})\ne(1,1,1)\), Step 3 gives a strictly larger norm term, and multiplying that strict inequality by \(\kappa>0\) preserves strictness.  Adding the equal linear parts to the strict norm comparison yields
\[
\phi_{r,\kappa}(1,1,1)<\phi_{r,\kappa}(x,y,z_{\mathrm{sp}}),
\]
so \((1,1,1)\) is the strict minimizer. % lean: frobenius_center_certificate
\end{proof}
